% !TEX root = ../risk_vs.tex
\section{Conclusion and Limitations}\label{sec:conclusion}

We show that the post-sample performance of published cross-sectional return predictors is remarkably similar to that of data-mined benchmarks. This result holds for most types of research we examine, including research that is risk-based or includes the support of a mathematical equilibrium model. Research that is agnostic about the theoretical foundation for predictability shows some signs of outperformance, but the magnitude is modest. The statistical implication is that whether a predictor is found in a journal or is data mined has little effect on mean inferences about post-sample performance (Equation \eqref{eq:painfully_clear}).

Beyond statistical inference, our findings suggest four deeper implications about cross-sectional stock return predictability: (1) empirical evidence is more informative than theoretical evidence for post-sample prediction, (2) investors do not learn about risk from academic research, (3) data mining is effective, and (4) mispricing is the primary driver. These implications come from analyzing the theoretical foundations of published predictors and their relationship with post-sample performance.

A limitation of our study is that we cannot identify the economic mechanism behind the lack of outperformance. The noisiness of predictor returns makes it difficult to determine the exact timing of decay, and thus makes it hard to separate publication effects (\citealt{mclean2016does}) from technological changes (\citealt{Chordia2014Have}). We illustrate this difficulty in the Internet Appendix \ref{sec:intapp-struct-break}.

A second limitation is that we study single-predictor strategies. For strategies that use many predictors, the factor structure and spanning are central questions. Table \ref{tab:dm-theme} suggests that data-mined accounting predictors are to a significant extent spanned by the ideas in the CZ dataset, but a more systematic investigation is needed. 

Last, our study is limited to the peer review process as characterized by the CZ dataset. This dataset is composed of papers that study cross-sectional return predictability, published between the years 1973 and 2016. Each literature has its own norms and practices, which evolve over time. The extent to which our findings generalize is an important question for future research.

\clearpage

